\documentclass{exam-zh}
\usepackage{siunitx}
\usepackage{draftwatermark} 

\SetWatermarkText{工试群版权所有，仅供学习，严禁商用}           % 设置水印内容
\SetWatermarkLightness{0.9}             % 设置水印透明度 0-1
\SetWatermarkScale{0.3}                   % 设置水印大小 0-1 


\examsetup{
  page/size=a4paper,
  paren/show-paren=true,
  paren/show-answer=true,
  fillin/show-answer=false,
  solution/show-solution=show-stay,
  solution/label-indentation=false,
  page/foot-content={试题第;页（共;页）}
}

\everymath{\displaystyle}

\title{上海交通大学2025年平台预分专业能力测试考试}

\subject{数学和逻辑}



\begin{document}

% \information{
%   姓名\underline{\hspace{6em}},
%   座位号\underline{\hspace{15em}}
% }
% \warning{（在此卷上答题无效）}

% \secret

\maketitle

本试卷共 60 题。全卷满分 100 分。考试用时 120 分钟。


\begin{notice}
  \item 上海交通大学工科试验班25群（QQ 1027107274）版权所有，禁止商用
  \item 制作人员：方天{\CJKfontspec{华文宋体}宬}（策划/整理/排版）、铁镐镐（整理）、陈翎玺（排版）、王浩宇（整理）、尹琳、unicorn、F丶忆笙、蔡炎彬Florence、黄浩原、王子越、青柠、wyzzzzzzzzz、罗嘉康(欢迎加入NIMO喵~)、AAA椒椒酱批发、白眉唠叨、浅水清、kumorin、初音ミク、genero、共潮生、煨鳐喂鹞、csix、雨绫汐、無漪雾、李敬泽、阿汐今天吃药了吗、白泽.、星宸大海、季姬击鸡记、凌琪鸯、户山香澄、青柠檬店、NIMO@水木年华、房奇喵、ymy、LoD（欢迎加入nimo喵！）、栗悟饭与龟波气功、澄闪、晓山瑞希（烤p的传承）、朱航Micropuppy、墨予、E＝mc\textasciicircum{}2、Sabrin
  \item 本卷是回忆版，疏漏之处请见谅
\end{notice}



\section{%
  数学部分：共 30 小题，1 到 10 题每小题 1 分，11 到 30 题每小题 2 分。
}

% Updated math section: removed explicit \\\ linebreaks and added \paren[] to each question stem

% 1
\begin{question}
已知集合$S$中元素的个数记为 $|S|$。若集合 $A$、$B$ 满足 $|A|=7$, $|A\cap B|=5$, $|A\cup B|=10$，则 $|B|=$ \paren[]
\begin{choices}
  \item $7$
  \item $8$
  \item $9$
  \item $10$
\end{choices}
\end{question}

% 2
\begin{question}
数据 $3,3,5,6$ 的中位数是 \paren[]
\begin{choices}
  \item $3$
  \item $4$
  \item $4.5$
  \item $5$
\end{choices}
\end{question}

% 3
\begin{question}
数列 $a_n$ 的首项为 $a_1$，公比为 $q$，$a_1>0$，且 $a_1$ 为数列的最大项，则 $q$ 的取值范围为 \paren[]
\begin{choices}
  \item $(-1,1)$
  \item $(-\infty,0)\cup(0,1)$
  \item $(0,1)$
  \item $[-1,0)\cup(0,1]$
\end{choices}
\end{question}

% 4
\begin{question}
平面直线 $2x+3y=1$ 的一个方向向量是 \paren[]
\begin{choices}
  \item $(3,-2)$
  \item $(2,-3)$
  \item $(3,2)$
  \item $(2,3)$
\end{choices}
\end{question}

% 5
\begin{question}
$\rho=\dfrac{4}{3-4\sin\theta}$ 代表的图形为 \paren[]
\begin{choices}
  \item 圆
  \item 椭圆
  \item 双曲线
  \item 抛物线
\end{choices}
\end{question}

% 6
\begin{question}
已知命题“单调函数不是周期函数”，下列说法正确的是 \paren[]
\begin{choices}
  \item 逆命题为“周期函数不是单调函数”
  \item 否命题为“单调函数是周期函数”
  \item 逆否命题为“周期函数是单调函数”
  \item 以上全错
\end{choices}
\end{question}

% 7
\begin{question}
对于实数域$\mathbb R$,若集合$X$满足:对于任意非空开区间,都存在$X$的元素,则称$X$在实数域$\mathbb R$稠密,那么在对于$\mathbb R$,$X$在$\mathbb R$中不稠密的充分必要条件是 \paren[]
\begin{choices}
  \item 对于任意非空开区间，都不存在 $X$ 中的元素
  \item 存在非空开区间，不存在 $X$ 中的元素
  \item 对于任意非空开区间都存在 $X$ 的补集中的元素
  \item 存在非空开区，存在 $X$ 的补集中的元素
\end{choices}
\end{question}

% 8
\begin{question}
记 $10^{x/2}$ 的图像为 $C$，$C$ 与 $C_1$ 关于 $x=1$ 对称，$C_2$ 与 $C_1$ 关于 $y=x$ 对称，则 $C_2$ 对应的函数是 \paren[]
\begin{choices}
  \item $1-2\lg x$
  \item $2-2\lg x$
  \item $2\lg x+1$
  \item $2\lg x+2$
\end{choices}
\end{question}

% 9
\begin{question}
设 $a_n$ 表示 $7^n$ 除以 $100$ 的余数，求 $\sum_{i=1}^{2025} a_i$ \paren[]
\begin{choices}
  \item 50607
  \item 50600
  \item 50635
  \item 50699
\end{choices}
\end{question}

% 10
\begin{question}
体积为 1 的立方体 $ABCD-A'B'C'D'$ 中，切去 $A'-BDA$, $B'-ACB$, $C'-BDC$, $D'-CAD$（重复部分不计）后，剩下的几何体体积是？ \paren[]
\begin{choices}
  \item $\frac{1}{3}$
  \item $\frac{3}{5}$
  \item $\frac{1}{2}$
  \item $\frac{2}{5}$
\end{choices}
\end{question}

% 11
\begin{question}
函数 $f(x)=\dfrac{\ln^2 x}{x}$ 的极值情况为 \paren[]
\begin{choices}
  \item 仅有极小值
  \item 仅有极大值
  \item 有极大值和极小值
  \item 无极大值和极小值
\end{choices}
\end{question}

% 12
\begin{question}
已知复数 $z$ 满足 $|z|-z=2-4i$，求 $\frac{1}{z}=$ \paren[]
\begin{choices}
  \item $\dfrac{3}{25}+\dfrac{4}{25}i$
  \item $\dfrac{3}{25}-\dfrac{4}{25}i$
  \item $\dfrac{3}{5}+\dfrac{4}{5}i$
  \item $3-4i$
\end{choices}
\end{question}

% 13
\begin{question}
点 $A(-3,1)$，圆 $x^2+y^2=1$ 与 $x^2-2x+y^2+2y=0$ 相交于 $B,C$，求 $\triangle ABC$ 的外接圆半径 $R$ \paren[]
\begin{choices}
  \item $1$
  \item $2$
  \item $3$
  \item $4$
\end{choices}
\end{question}

% 14
\begin{question}
若 $a,b,c>0$ 且 $a+\dfrac{1}{b}=b+\dfrac{1}{c}=c+\dfrac{1}{a}$，则该等式是 $a=b=c$ 的 \paren[]
\begin{choices}
  \item 充分不必要条件
  \item 必要不充分条件
  \item 既不充分也不必要条件
  \item 充要条件
\end{choices}
\end{question}

% 15
\begin{question}
某个大学组织了一个暑期实践活动，一共有12位学生参加，其中有3位理科生和9位工科生。现将12人分为3组，每组4个人，那么，3位理科生同时分到一个组的概率是 \paren[]
\begin{choices}
  \item $\dfrac{1}{4}$
  \item $\dfrac{1}{11}$
  \item $\dfrac{3}{55}$
  \item $\dfrac{1}{55}$
\end{choices}
\end{question}

% 16
\begin{question}
抛物线 $y^2=4x$ 上点 $A,B$ 满足 $\mathbf{AF}=3\mathbf{FB}$（$F$ 为焦点），求 $AB$ 中点到准线的距离为多少 \paren[]
\begin{choices}
  \item $\dfrac{5}{3}$
  \item $\dfrac{8}{3}$
  \item $\dfrac{11}{3}$
  \item $4$
\end{choices}
\end{question}

% 17
\begin{question}
已知 $P(B)=0.3$, $P(B|A)=0.375$, $P(B|\overline{A})=0.25$，求 $P(A|B)$ 为多少 \paren[]
\begin{choices}
  \item $0.2$
  \item $0.35$
  \item $0.5$
  \item $0.65$
\end{choices}
\end{question}

% 18
\begin{question}
已知 $0<y<x<\pi$ 且 $\cos x\cos y+\sin x\sin y=\dfrac{1}{2}$，$\sin(2x)+\sin(2y)=\dfrac{2}{3}$，求 $\tan(2x)$ 为多少 \paren[]
\begin{choices}
  \item $\dfrac{\sqrt{15}-2}{\sqrt{5}+2\sqrt{3}}$
  \item $\dfrac{\sqrt{15}-2}{\sqrt{5}-2\sqrt{3}}$
  \item $\dfrac{\sqrt{15}+2}{\sqrt{5}-2\sqrt{3}}$
  \item $\dfrac{\sqrt{15}+2}{\sqrt{5}+2\sqrt{3}}$
\end{choices}
\end{question}

% 19
\begin{question}
设 $D=\{(x,y)|0<y\le x^2\le1,\;0\le x\le1\}$，求该图形绕 $y$ 轴旋转一周所成实心体的体积为多少 \paren[]
\begin{choices}
  \item $\dfrac{\pi}{2}$
  \item $\dfrac{\pi}{3}$
  \item $\dfrac{\pi}{4}$
  \item $\dfrac{2\pi}{3}$
\end{choices}
\end{question}

% 20
\begin{question}
下列命题：① 存在四个互不相等且模相等的空间向量，两两的点积相等；② 对任意三个空间向量，总存在向量 $\vec{d}$ 使其和三者的点积均相等。 \paren[]
\begin{choices}
  \item ① 对 ② 不对
  \item ① 对 ② 对
  \item ① 不对 ② 不对
  \item ① 不对 ② 对
\end{choices}
\end{question}

% 21
\begin{question}
已知 $x_1=a,\;x_2=b,\;x_n=x_{n-1}-x_{n-2}\ (n\ge3)$，求 $x_{100}$ 与 $S_{100}$ \paren[]
\begin{choices}
  \item $x_{100}=-a\quad S_{100}=2b-a$
  \item $x_{100}=-a\quad S_{100}=b-a$
  \item $x_{100}=-b\quad S_{100}=2b-a$
  \item $x_{100}=-b\quad S_{100}=b-a$
\end{choices}
\end{question}

% 22
\begin{question}
$(x-\sqrt{2})^{2025}$ 中，取出所有 $x$ 的奇次幂项组成多项式 $S(x)$，求 $S(\sqrt{2})$ 为多少 \paren[]
\begin{choices}
  \item $2^{2035}\sqrt{2}$
  \item $2^{2036}\sqrt{2}$
  \item $2^{2037}\sqrt{2}$
  \item $2^{2038}\sqrt{2}$
\end{choices}
\end{question}

% 23
\begin{question}
甲乙两人抛硬币甲抛三次乙抛两次，硬币向上向下的概率相同，甲抛正面的次数比乙抛的正面的次数大的概率是 \paren[]
\begin{choices}
  \item $\dfrac{3}{5}$
  \item $\dfrac{7}{12}$
  \item $\dfrac{1}{2}$
  \item $\dfrac{11}{32}$
\end{choices}
\end{question}

% 24
\begin{question}
墙上有一幅画，观察者到其上、下端垂直距离分别为$a$和$b$，请问观察者与墙的水平距离为多少时，可以使画的上、下端对应视线夹角最大？ \paren[]
\begin{choices}
  \item $\dfrac{ab}{a^2+b^2}$
  \item $\dfrac{a+b}{2}$
  \item $\dfrac{ab}{a+b}$
  \item $\sqrt{ab}$
\end{choices}
\end{question}

% 25
\begin{question}
某地街道呈现东-西、南-北向的网格状，相邻街距都为1．两街道相交的点称为格点．若以互相垂直的两条街道为轴建立直角坐标系，现有下述格点$(-1,2)(3,1)(3,5)(-2,4)(4,4)(6,6)$为报刊零售点．请确定一个格点（除零售点外）为发行站，使6个零售点沿街道到发行站之间路程的和最短。 \paren[]
\begin{choices}
  \item $(-1,4)$
  \item $(3,3)$
  \item $(3,4)$
  \item $(4,3)$
\end{choices}
\end{question}

% 26
\begin{question}
$f(x)=(x+3)(x+4)(x+5)(x+6)$ 的最小值为 \paren[]
\begin{choices}
  \item $-1.5$
  \item $-1$
  \item $-0.5$
  \item $-0.25$
\end{choices}
\end{question}

% 27
\begin{question}
平面正六边形 $ABCDEF$，以$A$为起点，其余顶点为终点的向量记为$\mathbf{a}_1,\mathbf{a}_2,\cdots,\mathbf{a}_5$，以$D$为起点其余顶点为终点的向量记为$\mathbf{d}_1,\mathbf{d}_2,\cdots,\mathbf{d}_5$，$m$和$M$分别是$(\mathbf{a}_s+\mathbf{a}_t+\mathbf{a}_r)\cdot(\mathbf{d}_j+\mathbf{d}_i+\mathbf{d}_k)$的最小值和最大值，其中$s,t,r\in\{1,2,3,4,5\}$，$i,j,k\in\{1,2,3,4,5\}$ \paren[]
\begin{choices}
  \item $m<0,\;M<0$
  \item $m<0,\;M>0$
  \item $m=0,\;M>0$
  \item $m<0,\;M=0$
\end{choices}
\end{question}

% 28
\begin{question}
在体积为 1 的三棱锥 $A-BCD$ 的侧棱 $AB,AC,AD$ 上分别取点 $E,F,G$，且 $AE:EB=AF:FC=AG:GD=2:1$。记 $O$ 为三平面 $BCG,CDE,DBF$ 的交点，则 $O-BCD$ 的体积为多少 \paren[]
\begin{choices}
  \item $\dfrac{1}{9}$
  \item $\dfrac{1}{8}$
  \item $\dfrac{1}{7}$
  \item $\dfrac{1}{4}$
\end{choices}
\end{question}

% 29
\begin{question}
若一个细胞不分裂而死亡，且一个细胞分裂为两个细胞和死亡的概率均为$\dfrac{1}{2}$。现在有两个细胞，则2次分裂后还有细胞存活的概率为 \paren[]
\begin{choices}
  \item $\dfrac{39}{64}$
  \item $\dfrac{25}{64}$
  \item $\dfrac{29}{64}$
  \item $\dfrac{31}{64}$
\end{choices}
\end{question}

% 30
\begin{question}
将正 11 边形用对角线划分为 9 个三角形（对角线在内部两两不相交），关于锐角三角形的数量，下列哪项正确？ \paren[]
\begin{choices}
  \item 存在分法使所有三角形都不是锐角三角形
  \item 存在分法使恰有 2 个锐角三角形
  \item 存在分法使至少有 3 个锐角三角形
  \item 任何分法所分出的三角形都恰有 1 个锐角三角形
\end{choices}
\end{question}

% End of math section

\section{%
  逻辑部分：共 30 小题，31 到 40 题每小题 1 分，41 到 60 题每小题 2 分。
}

% Updated logic section: removed explicit \\\ linebreaks and added \paren[] to each question stem

% 31
\begin{question}
已知某会议室必须经过预约才能使用，今日该会议室的预约全部取消，则可以推出下列各项中的哪一项？ \paren[]
\begin{choices}
  \item 该会议室今日无人使用
  \item 该会议室被占用召开重要会议
  \item 预约系统发生故障
  \item 今日无人预约使用该会议室
\end{choices}
\end{question}

% 32
\begin{question}
某会议要求参会者穿正装出席（A），没有正装者穿深色便装出席（B），张某穿浅色便装出席，请问他违反了哪一项？ \paren[]
\begin{choices}
  \item A
  \item B
  \item A和B
  \item 没有违反
\end{choices}
\end{question}

% 33
\begin{question}
部长称投诉量连续三个月下降，说明服务改进取得成效。做出这样判断的前提是哪个？ \paren[]
\begin{choices}
  \item 投诉渠道畅通
  \item 投诉量是服务的唯一指标
  \item 没有其他影响投诉量因素
  \item 处理投诉内容的部门员工工作量变大
\end{choices}
\end{question}

% 34
\begin{question}
某省高考取消英语听力，农村考生平均分上升10\%。某些支持取消听力的人认为此举推动了教育公平，质疑该观点，应该要调查（）\paren[]
\begin{choices}
  \item 城市考生是否失去优势
  \item 取消听力对未来发展的影响
  \item 其他科目成绩是否受影响
  \item 农村地区教育资源是否存在劣势
\end{choices}
\end{question}

% 35
\begin{question}
以下证据能最好反驳“电子书能完全取代纸质书”的是哪个？ \paren[]
\begin{choices}
  \item 电子书更便宜
  \item 纸质书有收藏价值
  \item 电子书依赖电量
  \item 年轻人更喜欢电子书
\end{choices}
\end{question}

% 36
\begin{question}
某家饭店中，一桌人边用餐边谈生意。其中，一个哈尔滨人，两个北方人，一个广东人，两个人只做食品生意，三个人只做家电生意。如果以上介绍涉及餐桌上所有的人，那么这一桌最少可能是几个人？最多可能是几个人？ \paren[]
\begin{choices}
  \item 最少可能是3人，最多可能是8人
  \item 最少可能是5人，最多可能是8人
  \item 最少可能是5人，最多可能是9人
  \item 最少可能是3人，最多可能是9人
  \item 无法确定
\end{choices}
\end{question}

% 37
\begin{question}
凡金属都是导电的。铜是导电的，因此铜是金属。下面哪项与上述推理结构最相似？ \paren[]
\begin{choices}
  \item 所有的鸟都是卵生动物，蝙蝠不是卵生动物，所以，蝙蝠不是鸟。
  \item 所有的鸟都是卵生动物，天鹅是鸟，所以天鹅是卵生动物。
  \item 所有从事工商管理的都要学习企业管理，老陈是学习企业管理的，所以，老陈是从事工商管理工作的。
  \item 只有精通市场营销理论，才是一个合格的市场营销经理，老张精通市场营销理论，所以，老张一定是合格的市场营销经理。
  \item 华山险于黄山，黄山险于泰山，所以华山险于泰山。
\end{choices}
\end{question}

% 38
\begin{question}
大嘴鲈鱼只在有鲦鱼出现的河中且长有浮藻的水域里生活。漠亚河中没有大嘴鲈鱼。从上述断定能得出以下哪项结论? \paren[]

Ⅰ．鲦鱼只在长有浮藻的河中才能发现。

Ⅱ．漠亚河中既没有浮藻，又发现不了鲦鱼。

Ⅲ．如果在漠亚河中发现了鲦鱼，则其中肯定不会有浮藻。
\begin{choices}
  \item 只有Ⅰ。
  \item 只有Ⅱ。
  \item 只有Ⅲ。
  \item 只有Ⅰ和Ⅱ。
  \item Ⅰ、Ⅱ和Ⅲ都不是。
\end{choices}
\end{question}

% 39
\begin{question}
甲说：“除非所有的疾病都必然有确定的诱因，否则有些疾病可能难以预防。”乙说：“我不同意你的看法。”以下哪项断定，能准确表达乙的看法？ \paren[]
\begin{choices}
  \item 有些疾病必然没有确定的诱因，但所有疾病都可能加以预防
  \item 所有疾病都必然有确定的诱因，但有些疾病可能难以预防
  \item 有些疾病可能没有确定的诱因，但所有疾病都必然可以预防
  \item 所有疾病都可能没有确定的诱因，但有些疾病可能加以预防
  \item 有些疾病可能没有确定的诱因，但有些疾病可能加以预防
\end{choices}
\end{question}

% 40
\begin{question}
三个中国学生张林、赵强、李珊和三个外国留学生约翰、杰西、安娜暑假外出旅游，可供选择的旅游地有西安、杭州、大连和张家界。已经知道：

(1)每人只能去一个地方。

(2)凡是有中国学生去的地方，就必须有外国留学生去。

(3)凡是有外国留学生去的地方，就必须有中国学生去。

(4)约翰去西安或者杭州，赵强去张家界。

如果杰西去大连。则以下哪项一定为真? \paren[]

\begin{choices}
  \item 安娜去张家界
  \item 张林去大连
  \item 李珊去西安
  \item 约翰去杭州
  \item 张林和约翰
\end{choices}
\end{question}

% 41
\begin{question}
“冰淇淋销量越高，溺水率也越高，所以冰激凌销售导致溺水”对该推理说法正确的是 \paren[]
\begin{choices}
  \item 这是一个有效推理
  \item 这是一个无效推理，犯了诉诸恐惧的谬误
  \item 这是一个无效推理，犯了后以谬误（仅因为有先后顺序就得出因果关系）
  \item 这是一个无效推理，混淆了相关性和因果性
\end{choices}
\end{question}

% 42
\begin{question}
1.某企业减税5\%，研发投入增加3\%

2.某企业GDP增长，从百分之5.2降到百分之4.1

某人据此得出结论：减税并不能显著提高经济活力,他的论断问题在于 \paren[]
\begin{choices}
  \item 未定义经济活力的衡量标准
  \item 未比较减税前后企业研发投入
  \item 没有考虑其他因素对 GDP 增速的影响
  \item 将相关性等同于因果性
\end{choices}
\end{question}

% 43
\begin{question}
AI作品无需版权保护，因为其无人类创造性，以下哪项最能否定该观点 \paren[]
\begin{choices}
  \item 有些国家对 AI 作品提供版权保护
  \item AI 可能侵犯某些作家的版权
  \item AI 没有人的情感
  \item 人类需要设计算法和模型，因此前期过程有创造性
\end{choices}
\end{question}

% 44
\begin{question}
A 国 GDP 增长 8\%，通货膨胀 10\%；B 国 GDP 增长 3\%，通货膨胀 1\%。可以推出： \paren[]
\begin{choices}
  \item A 国经济发展更好
  \item B 国可能经济发展更好
  \item GDP 指标不可信
  \item 通货膨胀率与经济发展无关
\end{choices}
\end{question}

% 45
\begin{question}
专家说央行数字货币(CBDB)需要谨慎推行，以下最能支持专家观点的是： \paren[]
\begin{choices}
  \item 多国已试点推行CBDB
  \item 某国CBDB漏洞导致千万美金损失
  \item 现金使用率下降
  \item 私营加密货币波动性大
\end{choices}
\end{question}

% 46
\begin{question}
最近发现，19世纪80年代保存的海鸟标本的羽毛中，汞的含量仅为目前同一品种活鸟的羽毛汞含量的一半。由于海鸟羽毛中的汞的积累是海鸟吃鱼所导致的，这就表明现在海鱼中汞的含量比100多年前要高。以下哪项是上述论证的假设？\paren[]
\begin{choices}
  \item 进行羽毛汞含量检测的海鸟处于相同年龄段。
  \item 海鱼的汞含量取决于其活动海域的污染程度。
  \item 来源于鱼的汞被海鸟吸收后，残留在羽毛中的含量会随时间的变化而改变。
  \item 在海鸟的食物结构中，海鱼所占的比例，在19世纪80年代并不比现在高。
  \item 用于海鸟标本制作和保存的方法并没有显著减少海鸟羽毛中的汞含量。
\end{choices}
\end{question}

% 47
\begin{question}
张先生：应该向吸烟者征税，用以缓解医疗保健事业的投入不足。因为正是吸烟导致了许多严重的疾病。要吸烟者承担一部分费用来对付因他们不良习惯而造成的健康问题是完全合理的。

李女士：照您这么说，如果您经常吃奶油蛋糕，或者肥猪肉，也应该纳税。因为如同吸烟一样，经常食用高脂肪、高胆固醇的食物同样会导致许多严重的疾病。但是没有人会认为这样做是合理的。并且人们的危害健康的不良习惯数不胜数，都对此征税事实上无法操作。

以下哪项最为恰当地概括了张先生和李女士争论的焦点? \paren[]
\begin{choices}
  \item 张先生关于缓解医疗保健事业投入不足的建议是否合理?
  \item 有不良习惯的人是否应当对由此种习惯造成的社会后果负责?
  \item 食用高脂肪、高胆固醇的食物对健康造成的危害是否同吸烟一样严重？
  \item 由增加个人负担来缓解社会公共事业的投入不足是否合理?
  \item 通过征税的方式来纠正不良习惯是否合理?
\end{choices}
\end{question}

% 48
\begin{question}
某科研机构对市民所反映的一种奇异现象进行研究，该现象无法用已有的科学理论进行解释。助理研究员小王由此断言，该现象是错觉。

以下哪项如果为真,最可能使小王的断言不成立？ \paren[]
\begin{choices}
  \item 错觉都可以用已有的科学理论进行解释。
  \item 所有错觉都不能用已有的科学理论进行解释。
  \item 已有的科学理论尚不能完全解释错觉是如何形成的。
  \item 有些错觉不能用已有的科学理论进行解释。
  \item 有些错觉可以用已有的科学理论进行解释。
\end{choices}
\end{question}

% 49
\begin{question}
甲和乙任何一人都比丙、丁高。如果上述为真，再加上以下哪项，则可得出“戊比丁高”的结论？ \paren[]
\begin{choices}
  \item 戊比丙高
  \item 乙比甲高
  \item 乙比甲矮
  \item 戊比甲矮
  \item 戊比乙高
\end{choices}
\end{question}

% 50
\begin{question}
许多人不了解自己，也不设法去了解自己。这样的人可能想了解别人，但此种努力是要落空的，因为连自己都不了解的人不可能了解别人。由此可以得出结论:你要了解别人，就要了解自己。

以下哪项对上述论证的评价最为恰当？ \paren[]
\begin{choices}
  \item 上述论证所运用的推理是成立的。
  \item 上述论证有漏洞，因为它把得出某种结果的必要条件当做充分条件。
  \item 上述论证有漏洞，因为它不当地假设：每个人都可以了解自己。
  \item 上述论证有漏洞，因为它忽视了这种可能性：了解自己比了解别人更困难。
  \item 上述论证有漏洞，因为它基于个别性的事实轻率概括出一般性的结论。
\end{choices}
\end{question}

% 51
\begin{question}
观点：建造核电站促进碳达峰。

反对者的论据：

- 核电站建造耗时10年

- 光伏电站建设周期<2年

- 2030年要实现碳达峰

以下哪个是反对者的论证策略？ \paren[]
\begin{choices}
  \item 核电站具有危险性
  \item 核电站无法短期内减少碳排放
  \item 核电站成本效益小于光伏电站
  \item 清洁能源具有先进性
\end{choices}
\end{question}

% 52
\begin{question}
世卫组织警告，X病毒如果传播了，在低收入国家会因为医疗资源不均/匮乏导致死亡率高，可以推出什么结论？ \paren[]
\begin{choices}
  \item X病毒毒性更大
  \item X病毒在高收入国家死亡率可能更低
  \item X病毒在低收入国家已经蔓延
  \item 应该限制出境航班
\end{choices}
\end{question}

% 53
\begin{question}
环保组织称，除非禁止深海采矿，否则海洋生态会不可避免地遭到破坏。上述结论隐含了什么条件： \paren[]
\begin{choices}
  \item 深海开采已经导致生态遭到破坏
  \item 海洋生态效益大于经济效益（？
  \item 陆地矿物已被开采殆尽
  \item 现有开采技术无法避免生态破坏
\end{choices}
\end{question}

% 54
\begin{question}
研究发现，饮用纯净水的家庭比饮用自来水的家庭得肠胃炎的概率低20\%。有人据此认为纯净水比自来水更有利于健康。那么以下哪项如果为真，则既能解释数据也能反驳该论断？ \paren[]
\begin{choices}
  \item 饮用纯净水的家庭更注意饮食卫生
  \item 肠胃炎受病菌影响，与水质无关
  \item 自来水管道存在老化污染
  \item 该研究未区别桶装水与家用滤水器
\end{choices}
\end{question}

% 55
\begin{question}
（数独题）是个数独，挖了三个空让填： \paren[]
\begin{choices}
  \item 311
  \item 272
  \item 976
  \item 196
\end{choices}
\end{question}

% 56
\begin{question}
科学研究证明，非饱和脂肪酸含量高和饱和脂肪酸含量低的食物有利于预防心脏病。鱼通过食用浮游生物中的绿色植物使得体内含有丰富的非饱和脂肪酸“$\omega-3$”。而牛和其他反刍动物通过食用青草同样获得丰富的非饱和脂肪酸“$\omega-3$”。因此，多食用牛肉和多食用鱼对于预防心脏病都是有效的。以下哪项如果为真，最能削弱题干的论证？ \paren[]
\begin{choices}
  \item 在单位数量的牛肉和鱼肉中，前者非饱和脂肪酸“$\omega-3$”的含量要少于后者。
  \item 欧洲疯牛病的风波在全球范围内大大减少了牛肉的消费者，增加了鱼肉的消费者。
  \item 牛和其他反刍动物在反刍消化的过程中，把大量的非饱和脂肪酸转化为饱和脂肪酸。
  \item 实验证明，鱼肉中含有的非饱和脂肪酸“奥米加·3”比牛肉中含有的非饱和脂肪酸更易被人吸收。
  \item 统计表明，在欧洲内陆大量食用牛肉和奶制品的居民中患心脏病的比例，要高于在欧洲沿海大量食用鱼类的居民中患心脏病的比例。
\end{choices}
\end{question}

% 57
\begin{question}
在大型游乐园里，现场表演是刻意用来引导人群流动的。午餐时间的表演是为了减轻公园餐厅的压力；傍晚时间的表演则有一个完全不同的目的——鼓励参观者留下来吃晚餐。表面上不同时间的表演有不同的目的，但这背后，却有一个统一的潜在目标，即……

以下哪一项作为本段短文的结束语最为恰当?  \paren[]
\begin{choices}
  \item 尽可能地减少各游览点的排队人数。
  \item 吸引更多的人来看现场表演，以增加利润。
  \item 最大限度地避免由于游客出人公园而引起交通堵塞。
  \item 在尽可能多的时间里最大限度地发挥餐厅的作用。
  \item 尽可能地招揽顾客，希望他们再次来公园游览。
\end{choices}
\end{question}

% 58
\begin{question}
科学研究中使用的形式语言和日常生活中使用的自然语言有很大的不同，形式语言看起来像天书，远离大众，只有一些专业人士才能理解和运用；但其实这是一种误解，自然语言和形式语言的关系就像肉眼与显微镜的关系，肉眼的视域广阔，可以从整体上把握事物的信息；显微镜可以帮助人们看到事物的细节和精微之处，尽管用它看到的范围小。所以，形式语言和自然语言都是人们交流和理解信息的重要工具，把它们结合起来使用，具有强大的力量。

以下哪项如果为真，最能支持上述结论？ \paren[]
\begin{choices}
  \item 通过显微镜看到的内容肯能成为新的“风景”，说明形式语言可以丰富自然语言的表达，我们应重视形式语言。
  \item 正如显微镜下显示的信息最终还是要通过肉眼观察一样，形式语言表述的内容最终也要通过自然语言来实现，说明自然语言更基础。
  \item 科学理论如果仅用形式语言表达，很难被普通民众理解；同样，如果仅用自然语言表达，有可能变得冗长且很难表达准确。
  \item 科学的发展很大程度上改善了普通民众的日常生活，但人们并没有意识到科学表达的基础——形式语言的重要性。
  \item 采用哪种语言其实不重要，关键在于是否表达了真正想表达的思想内容。
\end{choices}
\end{question}

% 59
\begin{question}
某大学运动会即将召开，经管学院拟组建一支12人的代表队参赛，参赛队员将从该院4个年级学生中选拔，每个年级须在长跑、短跑、跳高、跳远、铅球等5个项目中选1~2项比赛，其余项目可任意选择。
一个年级如果选择长跑，就不能选短跑或跳高；一个年级如果选跳远，就不能选长跑或铅球，每名队员只参加一项比赛。

已知该院：

（1)每个年级均有队员被选拔进入代表队；

（2)每个年级被选拔进入代表队的人数各不相同；

（3)有两个年级的队员人数相乘等于另一个年级的队员人数。

如果某年级队员人数不是最少的，且选择长跑，那么对该年级来说，以下哪项不可能？\paren[]

\begin{choices}
  \item 选择铅球或跳远
  \item 选择短跑或铅球
  \item 选择短跑或跳远
  \item 选择长跑或跳高
  \item 选择铅球或跳高
\end{choices}
\end{question}

% 60
\begin{question}
第一个事实：

电视广告的效果越来越差。一项跟踪调查显示，在电视广告所推出的各种商品中，观众能够记住其品牌名称的商品的百分比逐年降低。

第二个事实：

在一段连续插播的电视广告中，观众印象较深的是第一个和最后一个，而中间播出的广告留给观众的印象，一般地说要浅得多。

以下哪项，如果为真，最能使得第二个事实成为对第一个事实的一个合理解释？ \paren[]
\begin{choices}
  \item 在从电视广告里见过的商品中，一般电视观众能记住其品牌名称的大约还不到一半。
  \item 近年来，被允许在电视节目中连续插播广告的平均时间逐渐缩短。
  \item 近年来，人们花在看电视上的平均时间逐渐缩短。
  \item 近年来，一段连续播出的电视广告所占用的平均时间逐渐增加。
  \item 近年来，一段连续播出的电视广告中所出现的广告的平均数量逐渐增加。
\end{choices}
\end{question}

\end{document}
